\documentclass[11pt,a4paper]{article}

\usepackage[margin=2.4cm,headheight=14pt]{geometry}
\usepackage{lmodern}
\usepackage[T1]{fontenc}
\usepackage{microtype}
\usepackage{booktabs}
\usepackage{tabularx}
\usepackage{array}
\usepackage{enumitem}
\usepackage{xcolor}
\usepackage{parskip}
\usepackage{tikz}
\usepackage{titlesec}
\usepackage{fancyhdr}
\usepackage[hidelinks]{hyperref}

% ---------------------------------------------------------------- palette ---
\definecolor{navy}{RGB}{13,27,42}       % deep night blue
\definecolor{navy2}{RGB}{27,38,59}      % lifted panel blue
\definecolor{emerald}{RGB}{16,185,129}  % bull green
\definecolor{gold}{RGB}{212,175,55}     % accent gold
\definecolor{mist}{RGB}{224,225,221}    % soft off-white

% ------------------------------------------------------------- typography ---
\newcolumntype{L}{>{\raggedright\arraybackslash}X}
\setlist{nosep}

\titleformat{\section}
  {\Large\bfseries\color{navy}}{\thesection}{0.8em}{}
  [{\color{emerald}\titlerule[1.4pt]}]
\titleformat{\subsection}
  {\large\bfseries\color{navy2}}{\thesubsection}{0.8em}{}

\pagestyle{fancy}
\fancyhf{}
\fancyhead[L]{\small\color{navy2}NafaIQ --- QA Test Plan}
\fancyhead[R]{\small\color{navy2}Team-H}
\fancyfoot[C]{\small\color{navy2}\thepage}
\renewcommand{\headrulewidth}{0.4pt}
\renewcommand{\headrule}{\color{emerald}\hrule height 0.9pt width \headwidth}

% cover stat block
\newcommand{\coverstat}[2]{%
  \begin{minipage}[t]{0.23\textwidth}\centering
    {\color{gold}\bfseries\LARGE #1}\\[3pt]
    {\color{mist}\footnotesize\scshape #2}
  \end{minipage}}

\begin{document}

% ============================================================== COVER PAGE ==
\begin{titlepage}
\begin{tikzpicture}[remember picture, overlay]
  % full-bleed background
  \fill[navy] (current page.south west) rectangle (current page.north east);

  % faint chart grid
  \begin{scope}[opacity=0.06]
    \foreach \y in {2,4,...,26}
      \draw[mist, line width=0.4pt]
        ([yshift=\y cm]current page.south west) -- ([yshift=\y cm]current page.south east);
    \foreach \x in {2,4,...,20}
      \draw[mist, line width=0.4pt]
        ([xshift=\x cm]current page.south west) -- ([xshift=\x cm]current page.north west);
  \end{scope}

  % rising market sparkline across the lower third
  \begin{scope}[shift={([yshift=6.2cm]current page.south west)}]
    \draw[emerald, opacity=0.55, line width=2.4pt, line join=round]
      (0,0.4) -- (2.6,1.5) -- (4.4,0.9) -- (7.0,2.6) -- (9.2,1.9)
      -- (12.0,3.6) -- (14.2,2.9) -- (17.4,4.8) -- (21.0,5.6);
    \fill[emerald, opacity=0.10]
      (0,0.4) -- (2.6,1.5) -- (4.4,0.9) -- (7.0,2.6) -- (9.2,1.9)
      -- (12.0,3.6) -- (14.2,2.9) -- (17.4,4.8) -- (21.0,5.6)
      -- (21.0,-6.2) -- (0,-6.2) -- cycle;
    \foreach \p in {(2.6,1.5),(7.0,2.6),(12.0,3.6),(17.4,4.8)}
      \fill[gold] \p circle (2.6pt);
  \end{scope}

  % top rule + team tag
  \draw[gold, line width=1.1pt]
    ([xshift=2.4cm, yshift=-2.6cm]current page.north west) --
    ([xshift=-2.4cm, yshift=-2.6cm]current page.north east);
  \node[anchor=north west, text=mist]
    at ([xshift=2.4cm, yshift=-3.2cm]current page.north west)
    {\large\scshape\textls[220]{Team-H \,\textcolor{gold}{\textbullet}\, Quality Assurance}};
  \node[anchor=north east, text=mist]
    at ([xshift=-2.4cm, yshift=-3.2cm]current page.north east)
    {\large\scshape\textls[220]{QA Week --- July 2026}};

  % title block (anchored off the page's LEFT edge, vertically near the middle)
  \node[anchor=west, text=white]
    at ([xshift=2.4cm, yshift=3.4cm]current page.west)
    {\fontsize{58}{58}\selectfont\bfseries NafaIQ};
  \node[anchor=west, text=emerald]
    at ([xshift=2.4cm, yshift=1.9cm]current page.west)
    {\fontsize{25}{25}\selectfont\bfseries Software Test Plan};
  \draw[gold, line width=2pt]
    ([xshift=2.4cm, yshift=1.0cm]current page.west) --
    ([xshift=8.6cm, yshift=1.0cm]current page.west);
  \node[anchor=north west, text=mist, text width=15.2cm]
    at ([xshift=2.4cm, yshift=0.5cm]current page.west)
    {\large Test strategy, design and execution for the NafaIQ investing and
     personal-finance platform --- web, mobile and API.};

  % stat strip
  \node[anchor=south, text=mist]
    at ([yshift=3.4cm]current page.south)
    {\begin{minipage}{0.86\paperwidth}\centering
       \coverstat{654}{test cases}\hfill
       \coverstat{228}{e2e journeys}\hfill
       \coverstat{1{,}757}{unit \& component tests}\hfill
       \coverstat{178}{api operations}
     \end{minipage}};

  % bottom rule + version
  \draw[gold, line width=1.1pt]
    ([xshift=2.4cm, yshift=2.2cm]current page.south west) --
    ([xshift=-2.4cm, yshift=2.2cm]current page.south east);
  \node[anchor=south, text=mist]
    at ([yshift=1.35cm]current page.south)
    {\small\scshape\textls[200]{Version 1.1 \,\textcolor{gold}{\textbullet}\, Jira Project KAN \,\textcolor{gold}{\textbullet}\, PSX / Supabase / FastAPI / React}};
\end{tikzpicture}
\end{titlepage}

% ================================================================ CONTENTS ==
\setcounter{page}{1}
\thispagestyle{empty}
{\hypersetup{linkcolor=navy}\tableofcontents}
\vfill
\begin{center}
\begin{tabularx}{\textwidth}{lL}
\toprule
\textbf{Product} & NafaIQ --- PSX investing and personal-finance platform (web, mobile, API) \\
\textbf{Prepared by} & Team-H (Quality Assurance) \\
\textbf{Test basis} & Live application (dev branch), OpenAPI schema (178 operations), product feature set \\
\textbf{Defect tracker} & Jira, project \texttt{KAN} \\
\textbf{Companion artefacts} & \texttt{NafaIQ\_Test\_Cases\_.xlsx} (test cases and results), \texttt{postman/} (API collections), \texttt{frontend/packages/e2e/} (automated journeys), \texttt{.github/workflows/} (CI gates) \\
\bottomrule
\end{tabularx}
\end{center}
\newpage

% ==================================================================== BODY ==
\section{Introduction}

\subsection{Purpose}
This document is the master test plan for the NafaIQ platform. It defines what
is tested, at which level, with which tools and data, in which environments,
and how results and defects are recorded. It is the first of the five QA-week
deliverables and the frame of reference for the other four:

\begin{enumerate}
  \item \textbf{Test plan} --- this document.
  \item \textbf{Test cases covering all features} --- \texttt{NafaIQ\_Test\_Cases\_.xlsx}:
        654 cases across 25 feature modules plus the automated-journey results sheet.
  \item \textbf{Automated end-to-end test cases} --- the Playwright suite
        (\texttt{@nafaiq/e2e}): 228 tests in 20 files across three browser projects,
        wired into CI and a nightly run, with a live failure-replay harness
        (Section~\ref{sec:e2e}).
  \item \textbf{Bugs reported in Jira} --- defects filed to project \texttt{KAN}
        under a fixed reporting convention (Section~\ref{sec:defects}).
  \item \textbf{API tasks} --- two Postman collections: a complete 178-operation
        API catalog and a 68-request scripted end-to-end QA suite
        (Section~\ref{sec:api}).
\end{enumerate}

\subsection{Product overview}
NafaIQ is a Pakistan-focused investing and personal-finance product covering
live PSX market data and analysis (signals V4), portfolio tracking, personal
finance (transactions, budgets, bills, goals, zakat), price/budget/goal alerts,
notifications, an AI assistant and AI reports, a learning hub with quizzes and
an AI tutor, and an admin subsystem. It ships as a React web app, an Expo
(React Native) mobile app, and a FastAPI backend on a Supabase (PostgreSQL)
database. The interface is bilingual (English and right-to-left Urdu).

\section{Test items and scope}

\subsection{Items under test}
\begin{itemize}
  \item \textbf{Web app} (\texttt{frontend/packages/web}) --- all user-facing routes,
        public and authenticated.
  \item \textbf{Backend API} (\texttt{backend/}) --- all 178 operations enumerated
        from the live OpenAPI schema, including the 35-endpoint admin subsystem.
  \item \textbf{Mobile app} (\texttt{frontend/packages/mobile}) --- screens, hooks
        and API clients at unit/component level.
  \item \textbf{Cross-cutting qualities} --- security and authorization, data
        integrity, localization/RTL, accessibility, PWA behaviour and resilience,
        performance and session handling, AI guardrails.
\end{itemize}

\subsection{Feature modules}
Test cases are organised into 25 modules, each on its own workbook sheet with a
stable ID prefix:

\begin{center}
\small
\begin{tabularx}{\textwidth}{llL}
\toprule
\textbf{Module} & \textbf{ID prefix} & \textbf{Examples of coverage} \\
\midrule
Authentication & TC-AUTH & sign-up/sign-in, validation, session, demo account \\
Dashboard & TC-DASH & greeting, KSE-100 strip, quick-add flows, AI recommendation \\
Portfolio & TC-PORT & holdings CRUD, cost averaging, sell paths, oversell guard \\
Finance & TC-FIN & transactions, budgets, bills, goals, summaries \\
Alerts / Notifications & TC-ALERT, TC-NOTIF & all condition types, channels, history \\
Watchlist & TC-WATCH & add/remove, duplicates, persistence \\
Market / Stock Detail & TC-MKT, TC-STK & index, movers, symbol pages \\
Market Data \& Macro & TC-MKTD & screener, sectors, macro data \\
Market Analysis V4 & TC-MA & signals endpoints, caching, degradation \\
AI Assistant / Reports & TC-AI, TC-REP & chat, streaming, quotas, report generation \\
Learn Hub & TC-LEARN & lessons, quizzes, AI tutor \\
User Profile & TC-PROF & settings, preferences, currency, theme \\
Admin Dashboard & TC-ADM & RBAC, users, roles, audit, flags \\
Security \& Authorization & TC-SEC & token handling, IDOR, credential scope \\
Data Integrity & TC-DATA & zakat correctness, arithmetic reconciliation, bulk delete \\
Cross-Module Flows & TC-FLOW & journeys that span portfolio, finance and alerts \\
Localization \& RTL & TC-L10N & Urdu interface, RTL layout, translations \\
Accessibility & TC-A11Y & WCAG checks, keyboard, labels, skip links \\
PWA \& Resilience & TC-PWA & offline behaviour, install, recovery \\
Performance \& Session & TC-PERF, TC-SESS & latency, caching, session expiry \\
AI Guardrails & TC-AIG & write confirmation, cross-user leakage, prompt abuse \\
Monetary Desk & TC-MON & monetary/macro desk views \\
\bottomrule
\end{tabularx}
\end{center}

\subsection{Out of scope}
\begin{itemize}
  \item Load and stress testing beyond the light concurrency probes recorded in
        Performance \& Session.
  \item Native mobile end-to-end automation on physical devices (mobile is
        covered at unit/component level plus a Pixel~7 mobile-viewport browser
        project in the E2E suite).
  \item Third-party systems themselves (Supabase, PSX data provider, AI model
        providers, e-mail delivery) --- only NafaIQ's integration with them.
  \item Subjective quality of AI-generated prose (factual/guardrail behaviour
        \emph{is} in scope under AI Guardrails).
\end{itemize}

\section{Test strategy overview}

Testing is layered so that each class of defect is caught by the cheapest layer
able to catch it, and every layer is automated where determinism allows.
Sections~\ref{sec:unit}--\ref{sec:manual} define each layer in detail.

\begin{center}
\small
\begin{tabularx}{\textwidth}{lllL}
\toprule
\textbf{Layer} & \textbf{Tooling} & \textbf{Volume} & \textbf{Primary target} \\
\midrule
Unit \& component & pytest, Vitest, Jest & 1{,}757 tests & logic, components, hooks, stores \\
Integration \& contract & pytest overrides, MSW & within the above & wiring, error taxonomy, schema drift \\
End-to-end (automated) & Playwright & 228 executions & real user journeys in a real browser \\
API (manual + scripted) & Postman & 178 + 68 requests & full API surface, hard journeys \\
Manual \& adversarial & live probes & 319 executed cases & security, data integrity, judgement calls \\
\bottomrule
\end{tabularx}
\end{center}

Every defect fix is re-verified at the level that found it, and the full
backend suite plus the Playwright suite are re-run after each fix wave. CI
(Section~\ref{sec:ci}) makes this automatic for every subsequent change.

\section{Unit and component testing}\label{sec:unit}

\begin{itemize}
  \item \textbf{Backend} --- pytest: \textbf{1457 passed, 4 skipped} at plan time,
        across 144 test files. API routes are exercised offline through
        dependency overrides and fake repositories; live-database tests
        auto-skip when credentials are absent so CI never depends on Supabase.
  \item \textbf{Web} --- Vitest with Testing Library and MSW: \textbf{179 tests in
        18 files} covering components, hooks, Redux slices and API contracts.
  \item \textbf{Mobile} --- Jest (jest-expo) with React Native Testing Library:
        \textbf{121 tests in 19 suites}; every screen test also asserts
        accessibility roles/labels per the mobile conventions.
\end{itemize}

Coverage is risk-ranked rather than exhaustive: money paths (portfolio
arithmetic, finance summaries, zakat), authentication and authorization, alert
evaluation, and the AI quota/refresh invariants take precedence over cosmetic
rendering.

\section{Integration and contract testing}\label{sec:integration}

Backend route$\rightarrow$service$\rightarrow$repository wiring is tested with the same
offline override pattern used at unit level, extended to whole-request flows:

\begin{itemize}
  \item \textbf{Cross-user isolation} --- every owned resource is probed with
        another user's identity; the expected result is a scoped 404, never a
        403 that confirms existence.
  \item \textbf{Error taxonomy} --- 401/403/404/422/429 contracts are asserted per
        route family, including the exact 401 messages the client refresh flow
        distinguishes.
  \item \textbf{Pagination and boundary edges} --- empty datasets, last-page
        behaviour, inclusive boundaries (e.g.\ the zakat nisab threshold).
  \item \textbf{Frontend contracts} --- web and mobile API mocks are generated
        from the real backend response schemas, so a backend field rename fails
        a frontend test rather than production. This guard exists because the
        finance-summary contract drifted once before QA week.
  \item \textbf{Policy middleware} --- CORS is pinned by test: allowed production
        and loopback origins, and anchored-regex rejection of
        \texttt{localhost.evil.com}-style suffix attacks.
\end{itemize}

\section{Automated end-to-end testing (Playwright)}\label{sec:e2e}

The \texttt{@nafaiq/e2e} package runs \textbf{228 tests in 20 files} against a real
browser and the real application stack (Vite dev server plus a local uvicorn
backend), across three projects:

\begin{center}
\begin{tabular}{lll}
\toprule
\textbf{Project} & \textbf{Browser} & \textbf{User} \\
\midrule
public-chromium & Desktop Chrome & Anonymous \\
authed-chromium & Desktop Chrome & Demo user \\
authed-mobile & Pixel 7 viewport & Demo user \\
\bottomrule
\end{tabular}
\end{center}

\subsection{Design rules}
All of the following are load-bearing for determinism:
\begin{itemize}
  \item \textbf{Journeys only} --- every spec is a user journey (open, act,
        observe on the destination page), asserted with web-first matchers on
        structure rather than polled market values.
  \item \textbf{One real login} --- a setup project signs in once with the demo
        account and saves storage state; only the Supabase session keys are
        persisted so no spec leaks app state into the next.
  \item \textbf{Deterministic AI} --- AI, assistant, tutor and report endpoints
        are stubbed with canned responses in the exact streaming (SSE) framing
        the client parses.
  \item \textbf{Console guard} --- any unexpected browser console error fails the
        test; known product bugs are allow-listed with a
        remove-when-fixed note tied to their Jira key.
  \item \textbf{Auth-settle guard} --- no interaction starts until the
        auth-resolution overlay has cleared, eliminating the swallowed-tap class
        of mobile flake (root cause of an early Pixel~7 defect).
  \item \textbf{Demo-account safety} --- demo-user writes live entirely in local
        state; browser journeys therefore never mutate shared backend data.
\end{itemize}

\subsection{Live failure replay}
A hard failure in the automated suite can be reproduced \emph{in front of the
tester} without manual setup. \texttt{pnpm test:e2e:live} runs the full suite and,
on the first hard failure, launches a headed Chromium, signs in through the
visible NafaIQ form using the stored demo credentials, and re-executes the
failing journey:

\begin{itemize}
  \item If the failure \textbf{reproduces}, Playwright pauses indefinitely with the
        browser and Inspector open at the failure state for live debugging.
  \item If it \textbf{passes} (a flake), the end state stays visible for five
        seconds and the browser closes.
\end{itemize}

\texttt{pnpm test:e2e:replay} replays any chosen failure from the latest JSON
report by index or by fingerprint. Replay mode disables traces and uses only
the console reporter so the demo password never lands in an artefact, and the
original suite report is not overwritten.

\section{API testing (Postman)}\label{sec:api}

Two collections, kept in \texttt{postman/} with Local and Dev environments:

\begin{itemize}
  \item \textbf{NafaIQ API --- Complete Catalog}: all \textbf{178} FastAPI operations,
        generated from the live OpenAPI schema and organised by the 27 API
        areas (finance, portfolio, signals-v4, admin, \ldots). Every request
        carries an assertion script; write requests are prefixed
        \emph{REVIEW BEFORE SEND}. This is the discovery and manual-inspection
        surface.
  \item \textbf{NafaIQ API --- End-to-End QA}: \textbf{68} scripted requests in nine
        ordered journey folders (setup/authentication, public market, portfolio,
        watchlist, personal finance, goals/budgets/bills, alerts and
        notifications, negative and authorization tests, optional AI).
        Requests chain dynamically captured IDs, assert on status and body, and
        clean up created records where the API supports it, so the collection
        is re-runnable and presentation-safe.
\end{itemize}

The catalog is also the coverage yardstick: it exposed that only ${\sim}40$ of 155
user-facing endpoints had cases before the expansion pass, which drove the
creation of the 13 newest workbook modules.

\section{Manual, exploratory and adversarial testing}\label{sec:manual}

Cases that automation cannot decide are executed manually against a live local
backend with real JWTs, with observed evidence recorded in the workbook's
\emph{Actual Result} column. The adversarial pass focused on:
security and authorization (RBAC on all admin endpoints, IDOR across two real
accounts, credential-scope confusion, token error hygiene), data integrity
(zakat nisab authority, clamping and boundaries, portfolio arithmetic
reconciliation, cost-averaging), input validation (impossible dates, negative
quantities, oversell), and AI guardrails (write-confirmation, cross-user data
requests, Roman-Urdu queries). Destructive probes were pointed at disposable
records or deliberately blocked (e.g.\ the untested bulk-delete endpoint,
TC-DATA-20, was not fired at the shared demo account).

\section{Test design and documentation}

All manual and API test cases live in \texttt{NafaIQ\_Test\_Cases\_.xlsx}:

\begin{itemize}
  \item \textbf{One sheet per module}, uniform columns: ID, module, scenario,
        preconditions, steps, test data, expected result, actual result,
        status, priority.
  \item \textbf{Statuses}: Pass, Fail, Partial, Blocked. A \emph{Blocked} case is
        never left bare --- its Actual Result states exactly what it needs
        (device, admin account, market window, failure injection, seeded data,
        human judgement). Nothing is marked Pass without recorded evidence.
  \item \textbf{Priorities}: P1 (core money paths, security, data integrity) through
        P4 (cosmetic). Current spread: 221 P1, 211 P2, 66 P3, 7 P4.
  \item \textbf{E2E Journey Results sheet}: one row per automated journey execution
        (149 rows = journeys $\times$ browser projects), carrying the automation
        source file, evidence path, run commit and timestamp.
  \item \textbf{Execution Summary sheet}: run metadata, defect focus, and the
        audit trail of each execution pass.
\end{itemize}

Traceability is bidirectional: automated cases name their spec file in the
workbook row, and defects name their test case in Jira.

\section{Estimation and prioritisation}\label{sec:estimation}

\subsection{Effort estimation}
QA week is a fixed five-day timebox with a two-person test team, giving a hard
budget of \textbf{10 person-days}. Effort was estimated bottom-up from the
module work-breakdown (Section~2.2) and allocated as follows:

\begin{center}
\small
\begin{tabularx}{\textwidth}{Llll}
\toprule
\textbf{Activity} & \textbf{Estimate} & \textbf{Share} & \textbf{Days} \\
\midrule
Test analysis and case design (workbook, 25 modules) & 2.5 pd & 25\% & Mon--Tue \\
Automation: E2E harness, then journey specs & 3.0 pd & 30\% & Mon--Wed \\
Test execution (manual + API waves, incl.\ re-tests) & 2.5 pd & 25\% & Wed--Fri \\
Environment, test identities and data setup & 0.5 pd & 5\% & Mon \\
Defect filing, fix verification and reporting & 1.0 pd & 10\% & continuous \\
Contingency (flake hunts, blocked-case triage) & 0.5 pd & 5\% & --- \\
\midrule
\textbf{Total} & \textbf{10 pd} & 100\% & 5 days \\
\bottomrule
\end{tabularx}
\end{center}

Estimation assumptions: the dev environment stays available all week; developer
capacity exists to fix Highest/High defects within the week (held --- 7 of 8
were fixed and re-verified); scope is frozen to the feature modules in
Section~2.2. Harness work is estimated separately from spec-writing because the
first automated journey carries the whole setup cost (authentication state,
deterministic stubs, anti-flake guards) while each subsequent journey is
marginal --- which is why automation was front-loaded into the first half of
the week.

Measured velocity, recorded for the next cycle's estimates: \textbf{319 live
case executions in the week} ($\approx$64 per day) alongside \textbf{228
automated journey executions} per suite run.

\subsection{Prioritisation under the five-day timebox}
A five-day box inverts the usual estimation question: not \emph{how long will
full coverage take}, but \emph{how much coverage fits, in what order}. The
answer is risk-first, encoded in every workbook row as a priority:

\begin{center}
\small
\begin{tabularx}{\textwidth}{clLl}
\toprule
\textbf{Priority} & \textbf{Cases} & \textbf{Definition} & \textbf{Timebox rule} \\
\midrule
P1 & 221 & Money paths, security, authorization, data integrity (zakat, portfolio arithmetic) & executed first, never cut \\
P2 & 211 & Core feature behaviour and validation & executed after P1 \\
P3 & 66 & Secondary flows, cosmetic states & as time allows \\
P4 & 7 & Nice-to-have polish & only if idle \\
\bottomrule
\end{tabularx}
\end{center}

Two consequences of this ordering are visible in the results: every filed
defect (Section~\ref{sec:defects}) came out of a P1/P2 case, and the cases that
ran out of runway are concentrated in the lower priorities and in the
resource-blocked bucket. The 233 Blocked cases --- each annotated with the
exact resource it needs (device, admin window, market hours, failure
injection) --- are deliberately retained as the costed input for the next
cycle's estimate rather than being dropped from scope.

\section{Test environment and data}

\begin{itemize}
  \item \textbf{Local stack} --- Vite dev server (web) + uvicorn (backend) +
        Supabase project; this is what Playwright boots and what manual API
        execution targeted (\texttt{127.0.0.1:8000}).
  \item \textbf{Deployed} --- backend on Railway, web on Vercel; the Postman Dev
        environment and an \texttt{E2E\_API\_URL} escape hatch point automation at it.
  \item \textbf{Test identities} --- the shared \textbf{demo account} (read-mostly by
        design; its writes stay client-side), a \textbf{second real account} for
        cross-user/IDOR checks, a purpose-made \textbf{QA super-admin} (role granted
        in the database, proving roles are DB-resolved rather than
        token-derived), and disposable targets for destructive operations. The
        demo account is deliberately \emph{not} an admin, because several cases
        assert it is refused.
  \item \textbf{Data hygiene} --- every server-side write in manual execution went
        to a throwaway record and was cleaned up; the demo baseline (net worth
        164{,}467 across 3 holdings) was re-verified identical after the pass.
\end{itemize}

\section{Continuous integration and nightly runs}\label{sec:ci}

\begin{itemize}
  \item \textbf{\texttt{ci.yml}} --- on every pull request and push: lint/typecheck,
        web unit, mobile unit, backend tests, and the two Chromium E2E
        projects, converging on a single required check (\texttt{ci-required})
        that branch protection and deploy gates key on.
  \item \textbf{\texttt{nightly-e2e.yml}} --- weekdays 08:00 PKT (before PSX opens):
        the full three-project matrix including the Pixel~7 mobile viewport,
        which is deliberately off the merge path. The nightly \textbf{always checks
        out the \texttt{dev} branch} (the integration branch), so a green nightly
        means anything promoted \texttt{dev}$\rightarrow$\texttt{main} is already
        QA-tested; a manual dispatch can override the ref.
  \item \textbf{Branch flow} --- feature branches $\rightarrow$ \texttt{dev}
        (integration) $\rightarrow$ \texttt{main} (production), with
        \texttt{ci-required} green as the promotion condition.
\end{itemize}

\section{Defect management}\label{sec:defects}

Defects are filed in Jira project \texttt{KAN} as they are found, under a fixed
convention:

\begin{itemize}
  \item \textbf{Title}: plain words stating what is broken and where.
  \item \textbf{Description}: exactly three parts --- \emph{Issue} (1--2 sentences),
        \emph{Steps} (short reproduction bullets), \emph{Expected} (1--2 sentences).
        Machine context (fingerprints, evidence paths, repro commands) goes in
        labels or a comment, never the description.
  \item \textbf{Severity} is judged by user impact and data/monetary risk;
        religiously sensitive computation (zakat) is automatically high.
  \item \textbf{Automated triage}: a runner parses the Playwright JSON report and
        files hard failures to \texttt{KAN} with a stable fingerprint label per
        failure, so re-occurrences comment on the existing ticket instead of
        duplicating it; infrastructure failures collapse to a single ticket,
        and a flood of new fingerprints halts for human confirmation.
\end{itemize}

Representative defects from this cycle (all found by planned cases, all but
one fixed and re-verified within the week):

\begin{center}
\small
\begin{tabularx}{\textwidth}{llLl}
\toprule
\textbf{Jira} & \textbf{Severity} & \textbf{Defect} & \textbf{State} \\
\midrule
KAN-52 & Highest & Zakat trusted a client-supplied nisab; \texttt{nisab=0} made PKR 1{,}000 of assets zakat-liable & Fixed \\
KAN-51 & High & Stock screener returned 500 on every request (service module missing) & Fixed \\
KAN-56 & High & Admin role history crashed for roles granted by a real admin & Fixed \\
KAN-53 & Medium & Add-holding accepted future and impossible purchase dates & Fixed \\
KAN-57 & Medium & Duplicate React keys on every lesson page (breadcrumb) & Fixed \\
KAN-58 & Medium & Missing skip-to-content link; unlabelled form fields (WCAG) & Fixed \\
KAN-55 & Medium & Market-analysis endpoint fully public --- gating is a product decision & Open \\
KAN-54 & Low & Auth 401 leaked Python codec internals to unauthenticated callers & Fixed \\
\bottomrule
\end{tabularx}
\end{center}

Verified-correct behaviour is recorded with the same rigour as defects (RBAC,
IDOR scoping, oversell protection, cost averaging, injection safety, zakat
clamping), so the presentation can state what is \emph{known good}, not only
what broke.

\section{Entry and exit criteria}

\textbf{Entry} (per execution pass): the build under test runs locally end to end;
test data baseline recorded; required identities available.

\textbf{Exit} (met at plan time):
\begin{itemize}
  \item Every workbook case carries a real result or a specific blocking reason
        --- no case left silently unexecuted.
  \item All automated suites green: backend 1457 passed, web 179, mobile 121,
        Playwright 228 executions across three projects (latest nightly:
        all three projects passing).
  \item No open Highest/High defect without a fix or an explicit product
        decision pending (sole remaining Fail: KAN-55, awaiting a gating
        decision).
  \item CI and nightly gates active on the integration branch.
\end{itemize}

\textbf{Suspension}: execution against shared data suspends immediately if a probe
mutates the shared demo baseline; the baseline is restored and verified before
resuming (exercised once during the pass, restored to the exact prior state).

\section{Risks and mitigations}

\begin{center}
\small
\begin{tabularx}{\textwidth}{LL}
\toprule
\textbf{Risk} & \textbf{Mitigation} \\
\midrule
Live PSX market data makes assertions time-of-day dependent & Journeys assert structure, not values; market-hours recorded fixtures planned for the still-blocked chart cases \\
Shared demo account could be polluted by test writes & Demo writes are client-side by design; server-side probes use throwaway records and a verified restore \\
AI output is non-deterministic & AI endpoints stubbed in E2E; guardrail behaviour (not prose quality) tested live \\
Flaky E2E failures erode trust in the gate & Retries in CI classify transients as flaky; live replay distinguishes real failures from flakes in one command; console guard keeps hidden errors visible \\
Concurrent merges can turn two green branches red together & The push-triggered CI run on the integration branch tip is the verdict of record; nightly re-checks \texttt{dev} daily \\
Blocked cases silently rot & Every Blocked row names its unblocking requirement, bucketed (device, admin, market window, failure injection, \ldots) for planned follow-up \\
\bottomrule
\end{tabularx}
\end{center}

\section{Deliverables and results snapshot}

\begin{center}
\small
\begin{tabularx}{\textwidth}{lLL}
\toprule
\textbf{Deliverable} & \textbf{Artefact} & \textbf{Status at plan time} \\
\midrule
Test plan & this document & Done \\
Test cases, all features & \texttt{NafaIQ\_Test\_Cases\_.xlsx}: 654 cases, 25 modules + journey results & 412 Pass / 8 Partial / 1 Fail / 233 Blocked-with-reason; 319 executed live this week \\
Automated E2E & \texttt{@nafaiq/e2e}: 228 tests, 20 files, 3 browser projects + live replay & Green locally, in CI and in the nightly (3/3 projects) \\
Bugs in Jira & project \texttt{KAN} & Defects filed throughout; this pass: 8 filed, 7 fixed and re-verified, 1 awaiting product decision \\
API tasks (Postman) & \texttt{postman/}: 178-operation catalog + 68-request scripted QA journeys, Local/Dev environments & Done; QA collection re-runnable with cleanup \\
\bottomrule
\end{tabularx}
\end{center}

Supporting (beyond the asked deliverables): 1757 unit/component/integration
tests across the three codebases, CI gating on every pull request, a nightly
cross-viewport run pinned to the integration branch, and an automated
Playwright-failure-to-Jira triage path.

\end{document}
